% FreeSerif is metric-compatible with Times New Roman, which the department
% template uses, and unlike most Times clones shipped with TeX it declares the
% Greek script, which polyglossia requires. FreeSans/FreeMono are its matching
% sans and monospace faces.
\usepackage{fontspec}

\setmainfont{FreeSerif.otf}[
  BoldFont       = FreeSerifBold.otf,
  ItalicFont     = FreeSerifItalic.otf,
  BoldItalicFont = FreeSerifBoldItalic.otf,
]

\setsansfont{FreeSans.otf}[
  BoldFont       = FreeSansBold.otf,
  ItalicFont     = FreeSansOblique.otf,
  BoldItalicFont = FreeSansBoldOblique.otf,
]

\setmonofont{FreeMono.otf}[
  BoldFont       = FreeMonoBold.otf,
  ItalicFont     = FreeMonoOblique.otf,
  BoldItalicFont = FreeMonoBoldOblique.otf,
  Scale          = MatchLowercase,
]
